% appendix/B_estimation.tex -- Appendix B (W1). B.1 Numerical solution (W1a); B.2 moments and weighting, B.3 assets and reentry fit (Figure B1),
% B.4 log-preference re-estimation (Table B1, Figure B2) added by W1b (BLUEPRINT 7, "Appendix B"; 12:155-210, 12:692-818). No visible placeholder text (Override O1).
% Vintage rule (BLUEPRINT Decision 8): Figure B1 model series are the 2026 files F:h800_meana_edit_c.tex and h800_mediana_edit_c.tex (footnoted);
% F:h800_reentry_edit.tex has one vintage. B.4 is the May 2025 robustness exercise (T:estimated_parameters_lnc.tex, F:*lnc*).
% Provenance: 12:NN = frot/sub2025/12_app_quant.tex; 11:NN = 11_app_data.tex.
\section{Solution, Estimation, and Fit}\label{app:estimation}

\subsection{Numerical solution}\label{app:solution}

The problem \eqref{eq:bellman} is solved by backward induction with value function iteration, starting from the terminal age $T=85$, at which the continuation value is zero, and proceeding to age 30. % src: 12:132-135; 5:23
The continuous state variables are discretized. The grid has 30 points for assets, 10 for the wage, and 10 for average indexed earnings; together with the two claiming states, the three health states, and the two previous-participation states, this gives the finite state space on which the value function is computed directly. % src: 12:140-143
% The grid bounds appear only in commented text of the May 2025 source (12:140-142) and are not stated (blueprint 2.3, "Do not say").
The expectation over next period's wage innovation $\varepsilon_{t+1}$ in \eqref{eq:wage} is computed by five-node Gauss--Hermite quadrature \citep{Judd1998}; the expectations over health and survival are finite sums over the three health states with the probabilities \eqref{eq:health} and \eqref{eq:survival}. % src: 12:146
The value function at points off the grid is evaluated by linear interpolation inside the grid and linear extrapolation outside it. % src: 12:146

The solution consists of decision rules for consumption, hours, and claiming as functions of the state $X_t$ and the parameters; the participation rule is the indicator that the hours rule is positive, and next period's assets follow from the budget constraint \eqref{eq:budget}. % src: 12:137-138
Simulated histories are generated by forward induction: an initial state is drawn for each simulated individual, the age-30 decision rules are applied, the age-31 state is generated from the transition probabilities and the realized shocks, and the process is repeated through age $T$. % src: 12:149-152
Initial states are drawn from the empirical joint distribution of wages, health, assets, and labor force participation of individuals aged 28--32 in the PSID sample, with the mean of log wages adjusted to the estimated age-30 profile within each health group; initial average indexed earnings are imputed from five years of work at 2,000 hours per year at the age-30 wage, following \citet{french2005effects}. % src: 11:222
\autoref{app:data} reports the initial distribution.

\subsection{Moments and weighting}\label{app:moments}

The second step estimates $\mathcal{P}=(\gamma,\nu,\allowbreak\phi_n(h{=}0),\allowbreak\phi_n(h{=}1,2),\allowbreak\phi_{RE},L,\beta)$ by the method of simulated moments, given the first-step inputs $\chi$: the net return $r$, the ability process $(\rho,\sigma_\varepsilon^2)$, the wage profiles $W(h,t)$, the health transition probabilities, the survival probabilities, the medical spending function, and the pension, Social Security, disability insurance, and tax rules. % src: 12:156; 12:133
Poor health and disability share one participation cost because the PSID has few observations of individuals with a disability. % src: 12:156-158
The estimator \eqref{eq:msm} minimizes the weighted distance between the target profiles of Appendix~\ref{app:targets} and the corresponding profiles simulated from the model. % src: 12:159-165
With $\mathbb{T}=40$ ages, the moment vector is % src: 12:167-180
\begin{equation}\label{eq:moments}
\hat{\varphi}(\mathcal{P};\chi)=
\begin{bmatrix}
\E[p_{iht}\mid h,t]-\int p_t(X,\mathcal{P},\chi)\,dF_{h,t}(X)\\[0.3em]
\E[n_{iht}\mid n_{iht}>0,\,h,t]-\int n_t(X,\mathcal{P},\chi)\,dF_{h,t}(X\mid n_t>0)\\[0.3em]
\E[a_{it}\mid t]-\int a_t(X,\mathcal{P},\chi)\,dF_{t}(X)\\[0.3em]
\operatorname{med}[a_{it}\mid t]-\operatorname{med}_{F_t}[a_t(X,\mathcal{P},\chi)]\\[0.3em]
\E[re_{it}\mid t]-\int re_t(X,\mathcal{P},\chi)\,dF_{t}(X)
\end{bmatrix},
\end{equation}
% src: 12:169-180 (moment vector; the mean and median asset rows are written out separately here, and the asset and reentry rows are not conditioned on health, as stated at 12:190-192 and 11:283)
for $t\in\{30,\dots,69\}$ and $h\in\{\text{healthy},\,\text{unhealthy}\}$, where $p_t$, $n_t$, $a_t$, and $re_t$ are the participation, hours, asset, and reentry rules implied by the decision rules of Appendix~\ref{app:solution}, $F_t$ is the distribution of the state at age $t$ in the simulated cohort, and $F_{h,t}$ is that distribution conditional on health $h$; the hours rows condition on positive hours in the data and in the simulation, since the target is hours conditional on participation. % src: 12:196; 12:188 (hours conditional on participation)
The blocks are: labor force participation by health (healthy, unhealthy) at ages 30--69, $2\mathbb{T}$ moments; hours worked conditional on participation by health, $2\mathbb{T}$ moments; mean and median assets, $2\mathbb{T}$ moments; and the mean labor market reentry rate, $\mathbb{T}$ moments---$7\mathbb{T}=280$ moments in all, for seven parameters, which leaves 273 degrees of freedom for the $\chi^2$ statistic of \autoref{tab:est}. % src: 12:186-195; T:estimated_parameters.tex:21
The weighting matrix $\hat{\mathbf{W}}_I$ is the inverse of the $7\mathbb{T}\times 7\mathbb{T}$ diagonal matrix whose diagonal is that of the variance--covariance matrix of the data moments; $I$ is the number of individuals in the data and $\varsigma$ the ratio of the number of data observations to the number of simulated observations. % src: 12:167; 12:197
The identifying assumptions follow \citet{french2005effects}: work decisions do not affect health or wages, and age changes work and saving incentives but not preferences. % src: 12:199

The estimation proceeds in five steps. % src: 12:201-203
First, the empirical life-cycle profiles of participation, hours, assets, and reentry are estimated from the PSID (Appendix~\ref{app:targets}). Second, given the initial distribution of the state variables and $\chi$, matrices of health and wage shocks are drawn. Third, for a value of $\mathcal{P}$ the model is solved by value function iteration and the decision rules are used to simulate profiles of participation, hours, assets, and reentry. Fourth, the moment conditions \eqref{eq:moments} and the criterion in \eqref{eq:msm} are computed. Fifth, a new $\mathcal{P}$ is chosen and the third and fourth steps are repeated until the criterion is minimized. % src: 12:201-203

\subsection{Fit of assets and reentry}\label{app:fitassets}

\autoref{fig:fitassets} completes \autoref{fig:fit} with the three remaining moment blocks. % src: 12:214-224
The model's mean assets lie inside the 95 percent band through the 40s and below the data from the early 50s; the median is above the data from the 30s through the early 50s and below it from the mid-50s; at age 60 the model has 1.642 against 3.218 for the mean and 1.183 against 2.252 for the median, in hundreds of thousands of 2016 dollars, and at age 65 1.185 against 3.493 and 0.7582 against 2.385. % src: F:h800_meana_edit_c.tex rows 30-49, 50-69, 60, 65; F:h800_mediana_edit_c.tex rows 30-53, 54-69, 60, 65
The reentry rate is zero in the model at every age; in the data it lies between 1.73 and 4.47 percent, with the maximum at age 63. % src: F:h800_reentry_edit.tex (model rows all 0; data minimum 1.73287 at age 37, maximum 4.47199 at age 63)
\autoref{subsec:fit} relates the asset misfit to the absence of a bequest motive and to the saving-subsidy reforms.

\begin{figure}[!htbp]
\centering
\caption{Model Fit: Assets and Labor Market Reentry}
\label{fig:fitassets}
\vspace{6pt}
\begin{tabular}{cc}
\begin{tikzpicture}
\begin{axis}[%
width=5.706cm,
height=4cm,
at={(0cm,0cm)},
scale only axis,
xmin=30,
xmax=70,
xlabel style={font=\color{white!15!black}},
xlabel={Age},
ymin=-5,
ymax=10,
ylabel style={font=\color{white!15!black}},
ylabel={Assets (\$100{,}000s)},
axis background/.style={fill=white},
%title style={font=\bfseries},
title={(A) Mean Assets},
legend style={at={(axis cs:32,5)}, anchor=south west, legend cell align=left, align=left, draw=white!15!black, font=\scriptsize},
grid=both,
grid style={line width=0.2pt, draw=gray!30},
minor grid style={line width=0.1pt, draw=gray!10},
minor tick num=1
]

% Main line (lfprlr0) - blue solid line
\addplot [color=black, line width=1pt]
  table[row sep=crcr]{%, y expr=\thisrowno{1}*100]{%
30	0.44883396	\\
31	0.42199558	\\
32	0.54018803	\\
33	0.78561033	\\
34	0.71175356	\\
35	0.74541797	\\
36	0.85600565	\\
37	0.95088265	\\
38	1.0613716	\\
39	1.1398318	\\
40	1.2659439	\\
41	1.2972758	\\
42	1.3919758	\\
43	1.535846	\\
44	1.6013677	\\
45	1.7205047	\\
46	1.8064766	\\
47	1.9908018	\\
48	2.0085264	\\
49	2.0823645	\\
50	2.2407959	\\
51	2.3326437	\\
52	2.556798	\\
53	2.5850001	\\
54	2.7065999	\\
55	2.8947089	\\
56	2.9723681	\\
57	2.9223428	\\
58	3.1388551	\\
59	3.1621034	\\
60	3.2175482	\\
61	3.620027	\\
62	3.4022311	\\
63	3.50429	\\
64	3.4009493	\\
65	3.4932815	\\
66	3.5353909	\\
67	3.6457606	\\
68	3.1442149	\\
69	3.8660995	\\
};
\addlegendentry{Data}

% Upper bound (standard error high)
\addplot[name path=upper, draw=none, forget plot] 
  table[row sep=crcr]{%,y expr=\thisrowno{1}*100]{%
%age1	plus2se	
30	1.5034447	\\
31	1.3318214	\\
32	1.3825211	\\
33	1.5269781	\\
34	1.4072133	\\
35	1.3745624	\\
36	1.4530003	\\
37	1.5145966	\\
38	1.5909592	\\
39	1.6398946	\\
40	1.7404519	\\
41	1.742667	\\
42	1.8167697	\\
43	1.9404281	\\
44	1.9883948	\\
45	2.0936916	\\
46	2.1680532	\\
47	2.3425176	\\
48	2.3583419	\\
49	2.4231668	\\
50	2.5766613	\\
51	2.6659567	\\
52	2.8871246	\\
53	2.9148088	\\
54	3.0332248	\\
55	3.2218368	\\
56	3.2981602	\\
57	3.260209	\\
58	3.4790351	\\
59	3.5159971	\\
60	3.5670526	\\
61	3.9753495	\\
62	3.7715802	\\
63	3.9004755	\\
64	3.8095388	\\
65	3.92889	\\
66	3.964179	\\
67	4.080106	\\
68	3.5837879	\\
69	4.3190544	\\
};

% Lower bound (standard error low)
\addplot[name path=lower, draw=none, forget plot] 
  table[row sep=crcr]{%,y expr=\thisrowno{1}*100]{%
%age1	minus2se	
30	-0.60577681	\\
31	-0.48783214	\\
32	-0.3021432	\\
33	0.04424252	\\
34	0.01629566	\\
35	0.11627356	\\
36	0.25901103	\\
37	0.38716866	\\
38	0.53178403	\\
39	0.63977097	\\
40	0.79143593	\\
41	0.85188656	\\
42	0.96718204	\\
43	1.1312658	\\
44	1.2143407	\\
45	1.347316	\\
46	1.4449018	\\
47	1.639086	\\
48	1.6587109	\\
49	1.7415622	\\
50	1.9049305	\\
51	1.9993306	\\
52	2.2264713	\\
53	2.2551913	\\
54	2.379975	\\
55	2.567581	\\
56	2.6465741	\\
57	2.5844784	\\
58	2.798675	\\
59	2.8082117	\\
60	2.8680457	\\
61	3.2647044	\\
62	3.032882	\\
63	3.1081044	\\
64	2.9923581	\\
65	3.057673	\\
66	3.1066009	\\
67	3.2114152	\\
68	2.7046419	\\
69	3.4131445	\\
};

% Fill between the paths
\addplot[blue!30] fill between[of=upper and lower];
\addlegendentry{95\% C.I.}

\addplot [color=red, line width=1.5pt, dashed]
  table[row sep=crcr]{%,y expr=\thisrowno{1}*100]{%
%age1	means	
30	0.528489028	\\
31	0.633230367	\\
32	0.739050697	\\
33	0.849863772	\\
34	0.960606134	\\
35	1.073980097	\\
36	1.183967902	\\
37	1.282999342	\\
38	1.377252345	\\
39	1.468645169	\\
40	1.559869506	\\
41	1.644474279	\\
42	1.722001208	\\
43	1.791600174	\\
44	1.860443712	\\
45	1.9207767	\\
46	1.976688014	\\
47	2.028959309	\\
48	2.072432524	\\
49	2.106475279	\\
50	2.131889566	\\
51	2.14796458	\\
52	2.151009084	\\
53	2.141524019	\\
54	2.116176358	\\
55	2.073609879	\\
56	2.020872452	\\
57	1.953956345	\\
58	1.870023323	\\
59	1.766523005	\\
60	1.641839629	\\
61	1.511149276	\\
62	1.357284975	\\
63	1.314517195	\\
64	1.255450127	\\
65	1.185440906	\\
66	1.189826815	\\
67	1.183263906	\\
68	1.165725954	\\
69	1.140764718	\\
  };
\addlegendentry{Model};

\end{axis}
\end{tikzpicture} & \begin{tikzpicture}
\begin{axis}[%
width=5.706cm,
height=4cm,
at={(0cm,0cm)},
scale only axis,
xmin=30,
xmax=70,
xlabel style={font=\color{white!15!black}},
xlabel={Age},
ymin=-5,
ymax=10,
ylabel style={font=\color{white!15!black}},
ylabel={Assets (\$100{,}000s)},
axis background/.style={fill=white},
%title style={font=\bfseries},
title={(B) Median Assets},
legend style={at={(axis cs:32,5)}, anchor=south west, legend cell align=left, align=left, draw=white!15!black, font=\scriptsize},
grid=both,
grid style={line width=0.2pt, draw=gray!30},
minor grid style={line width=0.1pt, draw=gray!10},
minor tick num=1
]

% Main line (lfprlr0) - blue solid line
\addplot [color=black, line width=1pt]
  table[row sep=crcr]{%, y expr=\thisrowno{1}*100]{%
30	-0.11895441	\\
31	-0.27639282	\\
32	0.0054679	\\
33	-0.05680846	\\
34	0.00530582	\\
35	0.05775673	\\
36	0.11935496	\\
37	0.283448	\\
38	0.29103227	\\
39	0.3317854	\\
40	0.3554958	\\
41	0.48654667	\\
42	0.49660314	\\
43	0.5230298	\\
44	0.75761317	\\
45	0.83516986	\\
46	0.95970209	\\
47	1.09713	\\
48	1.2602245	\\
49	1.2697649	\\
50	1.4709074	\\
51	1.4499263	\\
52	1.6115453	\\
53	1.7346132	\\
54	1.768639	\\
55	1.894712	\\
56	1.932978	\\
57	1.9691985	\\
58	2.1408199	\\
59	2.198014	\\
60	2.2522161	\\
61	2.329607	\\
62	2.4388831	\\
63	2.4980613	\\
64	2.5174607	\\
65	2.3850927	\\
66	2.403937	\\
67	2.3398796	\\
68	2.4810502	\\
69	2.4844838	\\
};
\addlegendentry{Data}

% Upper bound (standard error high)
\addplot[name path=upper, draw=none, forget plot] 
  table[row sep=crcr]{%,y expr=\thisrowno{1}*100]{%
30	0.31699318	\\
31	0.10011389	\\
32	0.35445438	\\
33	0.2506089	\\
34	0.29365724	\\
35	0.31921015	\\
36	0.36788288	\\
37	0.51851016	\\
38	0.51206232	\\
39	0.54050287	\\
40	0.55375253	\\
41	0.6729007	\\
42	0.67474134	\\
43	0.69266526	\\
44	0.9196122	\\
45	0.99146252	\\
46	1.1109516	\\
47	1.2443443	\\
48	1.4066153	\\
49	1.4124912	\\
50	1.6117372	\\
51	1.5897184	\\
52	1.74986	\\
53	1.8723206	\\
54	1.9050796	\\
55	2.0317431	\\
56	2.069126	\\
57	2.1103524	\\
58	2.2830432	\\
59	2.3452282	\\
60	2.3976214	\\
61	2.4770057	\\
62	2.5911442	\\
63	2.660554	\\
64	2.6842085	\\
65	2.563393	\\
66	2.5792043	\\
67	2.5160896	\\
68	2.66001	\\
69	2.6677378	\\
};

% Lower bound (standard error low)
\addplot[name path=lower, draw=none, forget plot] 
  table[row sep=crcr]{%,y expr=\thisrowno{1}*100]{%
30	-0.55490014	\\
31	-0.65290139	\\
32	-0.34352043	\\
33	-0.36422581	\\
34	-0.28304559	\\
35	-0.20369669	\\
36	-0.12917297	\\
37	0.04838584	\\
38	0.07000223	\\
39	0.12306792	\\
40	0.15724093	\\
41	0.30019264	\\
42	0.31846681	\\
43	0.35339433	\\
44	0.595616	\\
45	0.6788772	\\
46	0.80845444	\\
47	0.94991389	\\
48	1.1138318	\\
49	1.1270405	\\
50	1.3300777	\\
51	1.3101343	\\
52	1.4732306	\\
53	1.5969077	\\
54	1.6321967	\\
55	1.757679	\\
56	1.79683	\\
57	1.8280445	\\
58	1.9985966	\\
59	2.0507979	\\
60	2.1068127	\\
61	2.1822064	\\
62	2.2866238	\\
63	2.3355686	\\
64	2.3507148	\\
65	2.2067924	\\
66	2.2286696	\\
67	2.1636714	\\
68	2.3020905	\\
69	2.3012316	\\
};

% Fill between the paths
\addplot[blue!30] fill between[of=upper and lower];
\addlegendentry{95\% C.I.}

\addplot [color=red, line width=1.5pt, dashed]
  table[row sep=crcr]{%,y expr=\thisrowno{1}*100]{%
30	0	\\
31	0.169848043	\\
32	0.324047755	\\
33	0.458929333	\\
34	0.58588607	\\
35	0.708232142	\\
36	0.834945068	\\
37	0.953362006	\\
38	1.077344885	\\
39	1.195834593	\\
40	1.301419969	\\
41	1.400786892	\\
42	1.478278188	\\
43	1.5495248	\\
44	1.620944859	\\
45	1.681464662	\\
46	1.734871359	\\
47	1.775941666	\\
48	1.80198262	\\
49	1.82314822	\\
50	1.833265634	\\
51	1.827195632	\\
52	1.811864454	\\
53	1.782321958	\\
54	1.75352012	\\
55	1.695517106	\\
56	1.623180416	\\
57	1.535345572	\\
58	1.437292072	\\
59	1.328316533	\\
60	1.182584964	\\
61	1.045613071	\\
62	0.90142578	\\
63	0.871179625	\\
64	0.820927595	\\
65	0.758226595	\\
66	0.780832251	\\
67	0.794580994	\\
68	0.797297309	\\
69	0.785855602	\\
  };
\addlegendentry{Model};

\end{axis}
\end{tikzpicture}\\[1.5ex]
\multicolumn{2}{c}{\begin{tikzpicture}
\begin{axis}[%
width=5.706cm,
height=4cm,
at={(0cm,0cm)},
scale only axis,
xmin=30,
xmax=70,
xlabel style={font=\color{white!15!black}},
xlabel={Age},
ymin=0,
ymax=10,
ylabel style={font=\color{white!15!black}},
ylabel={Rate in \%},
axis background/.style={fill=white},
%title style={font=\bfseries},
title={(C) Mean Labor Market Reentry Rate},
legend style={at={(axis cs:31,9.7)}, anchor=north west, legend cell align=left, align=left, draw=white!15!black, font=\scriptsize},
grid=both,
grid style={line width=0.2pt, draw=gray!30},
minor grid style={line width=0.1pt, draw=gray!10},
minor tick num=1
]

% Main line (lfprlr0) - blue solid line
\addplot [color=black, line width=1pt]
  table[row sep=crcr]{%, y expr=\thisrowno{1}*100]{%
30		2.18517	\\
31		2.20335	\\
32		2.29117	\\
33		2.82388	\\
34		2.68548	\\
35		2.51887	\\
36		2.93026	\\
37		1.73287	\\
38		1.92231	\\
39		2.29121	\\
40		1.89049	\\
41		2.28087	\\
42		2.87458	\\
43		1.96387	\\
44		2.72082	\\
45		3.2679	\\
46		2.34941	\\
47		2.99802	\\
48		1.99764	\\
49		2.19266	\\
50		3.35012	\\
51		3.48878	\\
52		2.77504	\\
53		3.4707	\\
54		3.65526	\\
55		3.36258	\\
56		3.73983	\\
57		3.83127	\\
58		3.94909	\\
59		3.07238	\\
60		3.07621	\\
61		3.34096	\\
62		2.47159	\\
63		4.47199	\\
64		3.58776	\\
65		3.78918	\\
66		2.94066	\\
67		3.6062	\\
68		3.49552	\\
69		3.03359	\\
};
\addlegendentry{Data}

% Upper bound (standard error high)
\addplot[name path=upper, draw=none, forget plot] 
  table[row sep=crcr]{%,y expr=\thisrowno{1}*100]{%
30	3.52233	\\
31	3.48453	\\
32	3.49816	\\
33	3.97023	\\
34	3.77836	\\
35	3.57231	\\
36	3.94685	\\
37	2.7254	\\
38	2.89398	\\
39	3.24539	\\
40	3.38602	\\
41	3.69561	\\
42	4.20236	\\
43	3.22092	\\
44	3.9306	\\
45	4.43054	\\
46	3.47433	\\
47	4.08391	\\
48	3.0435	\\
49	3.21237	\\
50	4.36057	\\
51	4.48915	\\
52	3.77012	\\
53	4.46092	\\
54	4.63324	\\
55	4.33491	\\
56	4.7065	\\
57	4.81561	\\
58	4.93147	\\
59	4.07507	\\
60	4.06802	\\
61	4.34707	\\
62	3.48016	\\
63	5.51106	\\
64	4.64884	\\
65	4.87608	\\
66	4.04335	\\
67	4.7332	\\
68	4.6497	\\
69	4.25052	\\
};

% Lower bound (standard error low)
\addplot[name path=lower, draw=none, forget plot] 
  table[row sep=crcr]{%,y expr=\thisrowno{1}*100]{%
30	0.84801	\\
31	0.92216	\\
32	1.08417	\\
33	1.67754	\\
34	1.5926	\\
35	1.46544	\\
36	1.91367	\\
37	0.74035	\\
38	0.95064	\\
39	1.33704	\\
40	0.39496	\\
41	0.86613	\\
42	1.5468	\\
43	0.70682	\\
44	1.51103	\\
45	2.10526	\\
46	1.2245	\\
47	1.91214	\\
48	0.95177	\\
49	1.17295	\\
50	2.33967	\\
51	2.48841	\\
52	1.77995	\\
53	2.48048	\\
54	2.67729	\\
55	2.39025	\\
56	2.77315	\\
57	2.84692	\\
58	2.96672	\\
59	2.06969	\\
60	2.08439	\\
61	2.33486	\\
62	1.46303	\\
63	3.43292	\\
64	2.52668	\\
65	2.70228	\\
66	1.83796	\\
67	2.47919	\\
68	2.34135	\\
69	1.81667	\\
};

% Fill between the paths
\addplot[blue!30] fill between[of=upper and lower];
\addlegendentry{95\% C.I.}

\addplot [color=red, line width=1.5pt, dashed]
  table[row sep=crcr]{%
30	0	\\
31	0	\\
32	0	\\
33	0	\\
34	0	\\
35	0	\\
36	0	\\
37	0	\\
38	0	\\
39	0	\\
40	0	\\
41	0	\\
42	0	\\
43	0	\\
44	0	\\
45	0	\\
46	0	\\
47	0	\\
48	0	\\
49	0	\\
50	0	\\
51	0	\\
52	0	\\
53	0	\\
54	0	\\
55	0	\\
56	0	\\
57	0	\\
58	0	\\
59	0	\\
60	0	\\
61	0	\\
62	0	\\
63	0	\\
64	0	\\
65	0	\\
66	0	\\
67	0	\\
68	0	\\
69	0	\\
  };
\addlegendentry{Model};

\end{axis}
\end{tikzpicture}}
\end{tabular}
\begin{minipage}{0.95\textwidth}\footnotesize\setstretch{1.0}
\textit{Notes:} Black: PSID profiles from \eqref{eq:profile}; shaded: 95 percent confidence bands; red dashed: model. Ages 30--69. (A) Mean and (B) median nonpension assets in hundreds of thousands of 2016 dollars; (C) mean labor market reentry rate in percent, the share of individuals participating at age $t$ after not participating at $t-1$. Model series recomputed in 2026; data profiles and confidence bands are unchanged from the May 2025 version. Source: PSID 1968--2015.
% src: F:h800_meana_edit_c.tex, h800_mediana_edit_c.tex, h800_reentry_edit_paper.tex (series styles, titles, axis labels; the _paper copy changes the y-axis range of (C) only, figures/EDITS.md); 12:231-236 (May 2025 note); vintage sentence fixed in BLUEPRINT 9
\end{minipage}
\end{figure}

% \FloatBarrier removed (layout audit 10): the B.4 text may fill the page before the Figure B1 float page
\subsection{Re-estimation with logarithmic preferences}\label{app:lnc}

The model was also estimated with preferences that are separable in consumption and labor, % src: 12:694; 12:704
\begin{equation}\label{eq:lnc}
u(c_t,n_t;h_t)=\ln c_t-\varkappa\,\frac{n_t^{\,1+1/\epsilon}}{1+1/\epsilon}-\phi(h_t)\,\ind{P_t=1}-\phi_{RE}\,\ind{RE_t=1},
\end{equation}
% src: 12:705-707
where $\varkappa$ scales the disutility of hours, $\epsilon$ is the Frisch elasticity of labor supply, and $\phi(h_t)$ and $\phi_{RE}$ are fixed utility costs of participation and of reentry; the Frisch elasticity is set to $\epsilon=0.5$ following \citet{chetty2012bounds}, and the remaining parameters are estimated. % src: 12:710
The health and wage processes, the policy rules, the targets, and the method of simulated moments are those of the baseline. % src: 12:724
\autoref{tab:estlnc} reports the estimates: a discount factor of 0.97, a reentry cost of 36.15, participation costs of 0.53 for healthy and 0.59 for unhealthy workers, and $\varkappa=1.37\times 10^{-10}$; the costs are in units of flow utility and are not comparable in level with the time costs of \autoref{tab:est}. % src: T:estimated_parameters_lnc.tex:15-23
As in the baseline, the utility cost of working is higher for unhealthy than for healthy workers. % src: 12:728; T:estimated_parameters_lnc.tex:20-21
\autoref{fig:fitlnc} shows the fit. Mean assets are closer to the data at older ages than in the baseline (3.542 against 3.218 at age 60), while participation is matched less closely: healthy participation in the model is 96.86 percent against 83.94 in the data at age 50 and 68.44 against 64.34 at age 60. % src: F:h800_lnc_meana.tex row 60; F:h800_l0_lnc_edit.tex rows 50, 60; 12:729
The reform experiments of \autoref{sec:elimination} to \autoref{sec:other} were not re-run under this specification; \autoref{sec:scope} states what this implies for the sensitivity of the results.

\begin{table}[!htbp]
\centering
\caption{Estimated Parameters under Logarithmic Preferences}
\label{tab:estlnc}
\vspace{6pt}
\footnotesize
\begin{threeparttable}
% Copy of estimated_parameters_lnc.tex (Override O6). Edit: symbol $\kappa$ -> $\varkappa$ (row label only; distinguishes the hours-disutility scale from kappa_f of eq:mvpfcolumns; claims audit round 2, finding 5). Numbers untouched. Listed in tables/EDITS.md.

    
    \begin{tabular}{p{2.7cm} p{5.57cm} p{2.77cm} p{2.07cm}} % <-- Alignments: 1st column left, 2nd middle and 3rd right, with vertical lines in between
    \toprule
        \hline
Parameter & Definition & Estimates & S.E.\\[0.5ex]
 \hline





%$\gamma$& Consumption weight	& 0.5354 &\\%0.0042  \\
%$\nu$ & CRRA for flow utility 		& 1.7029 &\\% 0.0587\\
$\beta$ &Time discount factor		&
0.97     & 0.0015 \\
%L& Leisure endowment			& 4490.5  	&\\% 24.71\\
$\phi_{RE}$& Labor market reentry cost		& 
36.15  &1.2686\\
$\phi(h=0)$& Fixed cost of work, healthy 		& 0.53     &0.0082 \\
$\phi(h=1,2)$& Fixed cost of work, unhealthy 	& 0.59 		& 0.0152	\\
%$\varkappa$ & &1.37E-10 & 0.14E-10\\[1ex]
$\varkappa$ & Scaling factor for work disutility &$1.37\times 10^{-10}$ & $0.14\times 10^{-10}$\\[0.5ex]
%\multicolumn{2}{l}{$\chi^2$ statistic (\footnotesize{degrees of freedom} $= 273$)} & 1386.8\\[0.5ex]
 \hline
 \bottomrule
    \end{tabular}

	



\begin{tablenotes}[flushleft]\scriptsize
\item[] \textit{Notes:} Method of simulated moments with the utility function \eqref{eq:lnc}, the Frisch elasticity fixed at $\epsilon=0.5$, and the 280 moments of \autoref{tab:est}; standard errors in the last column. $\phi_{RE}$, $\phi(h{=}0)$, and $\phi(h{=}1,2)$ are in units of flow utility. May 2025 estimates.
% src: T:estimated_parameters_lnc.tex:15-23; 12:710; 12:742
\end{tablenotes}
\end{threeparttable}
\end{table}

\begin{figure}[!htbp]
\centering
\caption{Model Fit under Logarithmic Preferences}
\label{fig:fitlnc}
\vspace{6pt}
% Same layout as Figure 2 (long titles on the participation panels): zero column separation, centered in a \makebox of \textwidth.
\setlength{\tabcolsep}{0pt}%
\makebox[\textwidth][c]{\begin{tabular}{cc}
\begin{tikzpicture}
\begin{axis}[%
scale only axis,
width=5.706cm,
height=4cm,
at={(0cm,0cm)},
xmin=30,
xmax=70,
xlabel style={font=\color{white!15!black}},
xlabel={Age},
ymin=-5,
ymax=100,
ymajorgrids=true,
grid style=dashed,
ylabel style={font=\color{white!15!black}},
ylabel={Rate in \%},
axis background/.style={fill=white},
%title style={font=\bfseries},
title={(A) Healthy Labor Force Participation Rate},
legend style={at={(axis cs:32,0)}, anchor=south west, legend cell align=left, align=left, draw=white!15!black, font=\scriptsize},
grid=both,
grid style={line width=0.2pt, draw=gray!30},
minor grid style={line width=0.1pt, draw=gray!10},
minor tick num=1
]

% Main line (lfprlr0) - blue solid line
\addplot [color=black, line width=1pt]
  table[row sep=crcr, y expr=\thisrowno{1}*100]{%
30 0.8833822\\
31 0.8728638\\
32 0.8650124\\
33 0.8656822\\
34 0.8743696\\
35 0.873809\\
36 0.8829339\\
37 0.8818902\\
38 0.8847443\\
39 0.8934818\\
40 0.8900082\\
41 0.8684802\\
42 0.8632262\\
43 0.8659114\\
44 0.8597141\\
45 0.8664337\\
46 0.866869\\
47 0.8677746\\
48 0.8541223\\
49 0.8445035\\
50 0.8394285\\
51 0.8343548\\
52 0.8268231\\
53 0.8226036\\
54 0.8069361\\
55 0.7833442\\
56 0.7647346\\
57 0.7266455\\
58 0.7198906\\
59 0.6746376\\
60 0.6433743\\
61 0.5887213\\
62 0.4924323\\
63 0.4207888\\
64 0.3610537\\
65 0.2870897\\
66 0.2260264\\
67 0.21724\\
68 0.1962348\\
69 0.1569484\\
% 70 0.1681354\\
% 71 0.121358\\
% 72 0.1394618\\
% 73 0.1092736\\
% 74 0.1141775\\
% 75 0.0792269\\
% 76 0.0917666\\
% 77 0.0850947\\
% 78 0.0388536\\
% 79 0.0336195\\
% 80 0.058109\\
% 81 0.016314\\
% 82 0.0270162\\
% 83 -0.0089496\\
% 84 0.0236253\\
% 85 -0.0178846\\
};
\addlegendentry{Data}

% Upper bound (standard error high)
\addplot[name path=upper, draw=none, forget plot] 
  table[row sep=crcr,y expr=\thisrowno{1}*100]{%
30 0.9051051\\
31 0.8935855\\
32 0.8848711\\
33 0.88477\\
34 0.892873\\
35 0.8917789\\
36 0.9007229\\
37 0.8994025\\
38 0.9019169\\
39 0.9103903\\
40 0.9066504\\
41 0.8923349\\
42 0.8859496\\
43 0.8876875\\
44 0.8808528\\
45 0.8869151\\
46 0.88653\\
47 0.8869081\\
48 0.8729149\\
49 0.863117\\
50 0.858066\\
51 0.8529794\\
52 0.8455561\\
53 0.8413141\\
54 0.8257702\\
55 0.8021987\\
56 0.7835673\\
57 0.7460628\\
58 0.739598\\
59 0.6947576\\
60 0.663489\\
61 0.6094228\\
62 0.5136542\\
63 0.4426488\\
64 0.3833388\\
65 0.3098885\\
66 0.2493425\\
67 0.2413235\\
68 0.2207975\\
69 0.1832494\\
% 70 0.194256\\
% 71 0.1491448\\
% 72 0.1675596\\
% 73 0.1396055\\
% 74 0.1450964\\
% 75 0.1120739\\
% 76 0.1273876\\
% 77 0.1228318\\
% 78 0.0783176\\
% 79 0.0761687\\
% 80 0.1034711\\
% 81 0.0718914\\
% 82 0.080405\\
% 83 0.0504064\\
% 84 0.082339\\
% 85 0.0525297\\
};

% Lower bound (standard error low)
\addplot[name path=lower, draw=none, forget plot] 
  table[row sep=crcr,y expr=\thisrowno{1}*100]{%
30 0.8616592\\
31 0.852142\\
32 0.8451536\\
33 0.8465945\\
34 0.8558661\\
35 0.855839\\
36 0.8651446\\
37 0.8643779\\
38 0.8675715\\
39 0.8765733\\
40 0.8733659\\
41 0.8446254\\
42 0.8405027\\
43 0.8441353\\
44 0.8385755\\
45 0.8459523\\
46 0.8472079\\
47 0.848641\\
48 0.8353298\\
49 0.8258898\\
50 0.8207909\\
51 0.81573\\
52 0.80809\\
53 0.8038934\\
54 0.788102\\
55 0.7644898\\
56 0.7459021\\
57 0.7072282\\
58 0.7001832\\
59 0.6545175\\
60 0.6232595\\
61 0.5680198\\
62 0.4712104\\
63 0.3989288\\
64 0.3387687\\
65 0.2642909\\
66 0.2027102\\
67 0.1931565\\
68 0.1716722\\
69 0.1306474\\
% 70 0.1420148\\
% 71 0.0935712\\
% 72 0.111364\\
% 73 0.0789419\\
% 74 0.0832585\\
% 75 0.04638\\
% 76 0.0561455\\
% 77 0.0473576\\
% 78 -0.0006104\\
% 79 -0.0089297\\
% 80 0.012747\\
% 81 -0.0392634\\
% 82 -0.0263728\\
% 83 -0.0683057\\
% 84 -0.0350885\\
% 85 -0.0882989\\
};

% Fill between the paths
\addplot[blue!30] fill between[of=upper and lower];
\addlegendentry{95\% C.I.}

\addplot [color=red, line width=1.5pt, dashed]
  table[row sep=crcr,y expr=\thisrowno{1}*100]{%
age1	hlfprs	
30	0.968957203	\\
31	0.967528563	\\
32	0.968149038	\\
33	0.969477183	\\
34	0.969174977	\\
35	0.968055978	\\
36	0.96822487	\\
37	0.965898618	\\
38	0.965400062	\\
39	0.964252409	\\
40	0.962068966	\\
41	0.962097284	\\
42	0.962270133	\\
43	0.964990452	\\
44	0.963176433	\\
45	0.960663199	\\
46	0.963867188	\\
47	0.963036304	\\
48	0.961359094	\\
49	0.969137873	\\
50	0.968557758	\\
51	0.963202726	\\
52	0.960872576	\\
53	0.95022783	\\
54	0.934327299	\\
55	0.908797127	\\
56	0.881924198	\\
57	0.845669874	\\
58	0.798134328	\\
59	0.731264194	\\
60	0.684393064	\\
61	0.633671083	\\
62	0.423769508	\\
63	0.359203576	\\
64	0.308793456	\\
65	0.213959285	\\
66	0.197284684	\\
67	0.153846154	\\
68	0.141542507	\\
69	0.131132917	\\
  };
\addlegendentry{Model};

\end{axis}
\end{tikzpicture} & \begin{tikzpicture}
\begin{axis}[%
width=5.706cm,
height=4cm,
at={(0cm,0cm)},
scale only axis,
xmin=30,
xmax=70,
xlabel style={font=\color{white!15!black}},
xlabel={Age},
ymin=-5,
ymax=100,
ylabel style={font=\color{white!15!black}},
ylabel={Rate in \%},
axis background/.style={fill=white},
%title style={font=\bfseries},
title={(B) Unhealthy Labor Force Participation Rate},
legend style={at={(axis cs:32,0)}, anchor=south west, legend cell align=left, align=left, draw=white!15!black, font=\scriptsize},
grid=both,
grid style={line width=0.2pt, draw=gray!30},
minor grid style={line width=0.1pt, draw=gray!10},
minor tick num=1
]

% Main line (lfprlr0) - blue solid line
\addplot [color=black, line width=1pt]
  table[row sep=crcr,y expr=\thisrowno{1}*100]{%
30 0.8339695\\
31 0.830355\\
32 0.8042343\\
33 0.7983715\\
34 0.7995418\\
35 0.8378634\\
36 0.8354884\\
37 0.8686402\\
38 0.8411449\\
39 0.8375973\\
40 0.8386171\\
41 0.8234905\\
42 0.8157026\\
43 0.7707186\\
44 0.7771355\\
45 0.7791017\\
46 0.7758368\\
47 0.7928331\\
48 0.762918\\
49 0.710295\\
50 0.6906397\\
51 0.6569952\\
52 0.6095482\\
53 0.5680759\\
54 0.5703449\\
55 0.5509109\\
56 0.5310219\\
57 0.5219091\\
58 0.4718995\\
59 0.4247883\\
60 0.4079395\\
61 0.3693248\\
62 0.2990521\\
63 0.2519826\\
64 0.2448949\\
65 0.1944834\\
66 0.1670593\\
67 0.1475739\\
68 0.1256826\\
69 0.1083655\\
% 70 0.1681354\\
% 71 0.121358\\
% 72 0.1394618\\
% 73 0.1092736\\
% 74 0.1141775\\
% 75 0.0792269\\
% 76 0.0917666\\
% 77 0.0850947\\
% 78 0.0388536\\
% 79 0.0336195\\
% 80 0.058109\\
% 81 0.016314\\
% 82 0.0270162\\
% 83 -0.0089496\\
% 84 0.0236253\\
% 85 -0.0178846\\
};
\addlegendentry{Data}

% Upper bound (standard error high)
\addplot[name path=upper, draw=none, forget plot] 
  table[row sep=crcr,y expr=\thisrowno{1}*100]{%
30 0.9111052\\
31 0.9001306\\
32 0.8669667\\
33 0.8606798\\
34 0.860257\\
35 0.8955216\\
36 0.8874309\\
37 0.9166927\\
38 0.8892378\\
39 0.8830312\\
40 0.8826435\\
41 0.8928777\\
42 0.8792804\\
43 0.8314206\\
44 0.83197\\
45 0.8315616\\
46 0.8263692\\
47 0.8398106\\
48 0.8062539\\
49 0.7522238\\
50 0.7306355\\
51 0.6956045\\
52 0.6471571\\
53 0.6047442\\
54 0.6048974\\
55 0.5855613\\
56 0.5653728\\
57 0.5550599\\
58 0.5043322\\
59 0.4570315\\
60 0.4391904\\
61 0.4005467\\
62 0.3297285\\
63 0.2832072\\
64 0.2764793\\
65 0.2267636\\
66 0.1990062\\
67 0.1796762\\
68 0.1588332\\
69 0.142154\\
% 70 0.194256\\
% 71 0.1491448\\
% 72 0.1675596\\
% 73 0.1396055\\
% 74 0.1450964\\
% 75 0.1120739\\
% 76 0.1273876\\
% 77 0.1228318\\
% 78 0.0783176\\
% 79 0.0761687\\
% 80 0.1034711\\
% 81 0.0718914\\
% 82 0.080405\\
% 83 0.0504064\\
% 84 0.082339\\
% 85 0.0525297\\
};

% Lower bound (standard error low)
\addplot[name path=lower, draw=none, forget plot] 
  table[row sep=crcr,y expr=\thisrowno{1}*100]{%
30 0.7568338\\
31 0.7605796\\
32 0.7415019\\
33 0.7360631\\
34 0.7388266\\
35 0.7802051\\
36 0.7835459\\
37 0.8205878\\
38 0.7930518\\
39 0.7921634\\
40 0.7945906\\
41 0.7541034\\
42 0.7521245\\
43 0.7100167\\
44 0.7223011\\
45 0.7266418\\
46 0.7253045\\
47 0.7458558\\
48 0.719582\\
49 0.6683661\\
50 0.6506439\\
51 0.618386\\
52 0.5719393\\
53 0.5314076\\
54 0.5357924\\
55 0.5162606\\
56 0.496671\\
57 0.4887584\\
58 0.4394668\\
59 0.3925449\\
60 0.3766887\\
61 0.3381029\\
62 0.2683756\\
63 0.2207581\\
64 0.2133105\\
65 0.1622031\\
66 0.1351124\\
67 0.1154717\\
68 0.092532\\
69 0.0745769\\
% 70 0.1420148\\
% 71 0.0935712\\
% 72 0.111364\\
% 73 0.0789419\\
% 74 0.0832585\\
% 75 0.04638\\
% 76 0.0561455\\
% 77 0.0473576\\
% 78 -0.0006104\\
% 79 -0.0089297\\
% 80 0.012747\\
% 81 -0.0392634\\
% 82 -0.0263728\\
% 83 -0.0683057\\
% 84 -0.0350885\\
% 85 -0.0882989\\
};

% Fill between the paths
\addplot[blue!30] fill between[of=upper and lower];
\addlegendentry{95\% C.I.}

\addplot [color=red, line width=1.5pt, dashed]
  table[row sep=crcr,y expr=\thisrowno{1}*100]{%
30	0.802197802	\\
31	0.82183908	\\
32	0.854651163	\\
33	0.890052356	\\
34	0.910994764	\\
35	0.9342723	\\
36	0.933920705	\\
37	0.963265306	\\
38	0.969581749	\\
39	0.978798587	\\
40	0.977419355	\\
41	0.967065868	\\
42	0.953757225	\\
43	0.944134078	\\
44	0.949602122	\\
45	0.964622642	\\
46	0.974299065	\\
47	0.978723404	\\
48	0.975903614	\\
49	0.967244701	\\
50	0.933797909	\\
51	0.90619469	\\
52	0.87254902	\\
53	0.845440495	\\
54	0.803806735	\\
55	0.755244755	\\
56	0.69973545	\\
57	0.65037594	\\
58	0.591463415	\\
59	0.527972028	\\
60	0.489502762	\\
61	0.447537473	\\
62	0.376623377	\\
63	0.332050048	\\
64	0.290047393	\\
65	0.187557182	\\
66	0.153105862	\\
67	0.143564356	\\
68	0.118226601	\\
69	0.109697933	\\
  };
\addlegendentry{Model};

\end{axis}
\end{tikzpicture}\\[1.5ex]
\begin{tikzpicture}
\begin{axis}[%
scale only axis,
width=5.706cm,
height=4cm,
at={(0cm,0cm)},
xmin=30,
xmax=70,
xlabel style={font=\color{white!15!black}},
xlabel={Age},
ymin=-5,
ymax=10,
ymajorgrids=true,
grid style=dashed,
ylabel style={font=\color{white!15!black}},
ylabel={Assets (\$100{,}000s)},
axis background/.style={fill=white},
%title style={font=\bfseries},
title={(C) Mean Assets},
legend style={at={(axis cs:32,5)}, anchor=south west, legend cell align=left, align=left, draw=white!15!black, font=\scriptsize},
grid=both,
grid style={line width=0.2pt, draw=gray!30},
minor grid style={line width=0.1pt, draw=gray!10},
minor tick num=1
]

% Main line (lfprlr0) - blue solid line
\addplot [color=black, line width=1pt]
  table[row sep=crcr]{%, y expr=\thisrowno{1}*100]{%
30	0.44883396	\\
31	0.42199558	\\
32	0.54018803	\\
33	0.78561033	\\
34	0.71175356	\\
35	0.74541797	\\
36	0.85600565	\\
37	0.95088265	\\
38	1.0613716	\\
39	1.1398318	\\
40	1.2659439	\\
41	1.2972758	\\
42	1.3919758	\\
43	1.535846	\\
44	1.6013677	\\
45	1.7205047	\\
46	1.8064766	\\
47	1.9908018	\\
48	2.0085264	\\
49	2.0823645	\\
50	2.2407959	\\
51	2.3326437	\\
52	2.556798	\\
53	2.5850001	\\
54	2.7065999	\\
55	2.8947089	\\
56	2.9723681	\\
57	2.9223428	\\
58	3.1388551	\\
59	3.1621034	\\
60	3.2175482	\\
61	3.620027	\\
62	3.4022311	\\
63	3.50429	\\
64	3.4009493	\\
65	3.4932815	\\
66	3.5353909	\\
67	3.6457606	\\
68	3.1442149	\\
69	3.8660995	\\
};
\addlegendentry{Data}

% Upper bound (standard error high)
\addplot[name path=upper, draw=none, forget plot] 
  table[row sep=crcr]{%,y expr=\thisrowno{1}*100]{%
%age1	plus2se	
30	1.5034447	\\
31	1.3318214	\\
32	1.3825211	\\
33	1.5269781	\\
34	1.4072133	\\
35	1.3745624	\\
36	1.4530003	\\
37	1.5145966	\\
38	1.5909592	\\
39	1.6398946	\\
40	1.7404519	\\
41	1.742667	\\
42	1.8167697	\\
43	1.9404281	\\
44	1.9883948	\\
45	2.0936916	\\
46	2.1680532	\\
47	2.3425176	\\
48	2.3583419	\\
49	2.4231668	\\
50	2.5766613	\\
51	2.6659567	\\
52	2.8871246	\\
53	2.9148088	\\
54	3.0332248	\\
55	3.2218368	\\
56	3.2981602	\\
57	3.260209	\\
58	3.4790351	\\
59	3.5159971	\\
60	3.5670526	\\
61	3.9753495	\\
62	3.7715802	\\
63	3.9004755	\\
64	3.8095388	\\
65	3.92889	\\
66	3.964179	\\
67	4.080106	\\
68	3.5837879	\\
69	4.3190544	\\
};

% Lower bound (standard error low)
\addplot[name path=lower, draw=none, forget plot] 
  table[row sep=crcr]{%,y expr=\thisrowno{1}*100]{%
%age1	minus2se	
30	-0.60577681	\\
31	-0.48783214	\\
32	-0.3021432	\\
33	0.04424252	\\
34	0.01629566	\\
35	0.11627356	\\
36	0.25901103	\\
37	0.38716866	\\
38	0.53178403	\\
39	0.63977097	\\
40	0.79143593	\\
41	0.85188656	\\
42	0.96718204	\\
43	1.1312658	\\
44	1.2143407	\\
45	1.347316	\\
46	1.4449018	\\
47	1.639086	\\
48	1.6587109	\\
49	1.7415622	\\
50	1.9049305	\\
51	1.9993306	\\
52	2.2264713	\\
53	2.2551913	\\
54	2.379975	\\
55	2.567581	\\
56	2.6465741	\\
57	2.5844784	\\
58	2.798675	\\
59	2.8082117	\\
60	2.8680457	\\
61	3.2647044	\\
62	3.032882	\\
63	3.1081044	\\
64	2.9923581	\\
65	3.057673	\\
66	3.1066009	\\
67	3.2114152	\\
68	2.7046419	\\
69	3.4131445	\\
};

% Fill between the paths
\addplot[blue!30] fill between[of=upper and lower];
\addlegendentry{95\% C.I.}

\addplot [color=red, line width=1.5pt, dashed]
  table[row sep=crcr]{%,y expr=\thisrowno{1}*100]{%
%age1	means	
30		0.52848903	\\
31		0.63949028	\\
32		0.75641608	\\
33		0.88135533	\\
34		1.0146479	\\
35		1.15531	\\
36		1.2980178	\\
37		1.4347526	\\
38		1.5716783	\\
39		1.7089924	\\
40		1.8504481	\\
41		1.990505	\\
42		2.1279477	\\
43		2.2655075	\\
44		2.4060401	\\
45		2.543429	\\
46		2.6864321	\\
47		2.8260226	\\
48		2.9579886	\\
49		3.0826261	\\
50		3.2063248	\\
51		3.3196893	\\
52		3.418323	\\
53		3.503818	\\
54		3.572795	\\
55		3.6217814	\\
56		3.6555198	\\
57		3.6704038	\\
58		3.6589162	\\
59		3.6182893	\\
60		3.5423121	\\
61		3.4621445	\\
62		3.3563867	\\
63		3.3740218	\\
64		3.3360893	\\
65		3.2575333	\\
66		3.15897	\\
67		3.0522143	\\
68		2.9347336	\\
69		2.8075372	\\
  };
\addlegendentry{Model};

\end{axis}
\end{tikzpicture} & \begin{tikzpicture}
\begin{axis}[%
width=5.706cm,
height=4cm,
at={(0cm,0cm)},
scale only axis,
xmin=30,
xmax=70,
xlabel style={font=\color{white!15!black}},
xlabel={Age},
ymin=-5,
ymax=10,
ylabel style={font=\color{white!15!black}},
ylabel={Assets (\$100{,}000s)},
axis background/.style={fill=white},
%title style={font=\bfseries},
title={(D) Median Assets},
legend style={at={(axis cs:32,5)}, anchor=south west, legend cell align=left, align=left, draw=white!15!black, font=\scriptsize},
grid=both,
grid style={line width=0.2pt, draw=gray!30},
minor grid style={line width=0.1pt, draw=gray!10},
minor tick num=1
]

% Main line (lfprlr0) - blue solid line
\addplot [color=black, line width=1pt]
  table[row sep=crcr]{%, y expr=\thisrowno{1}*100]{%
30	-0.11895441	\\
31	-0.27639282	\\
32	0.0054679	\\
33	-0.05680846	\\
34	0.00530582	\\
35	0.05775673	\\
36	0.11935496	\\
37	0.283448	\\
38	0.29103227	\\
39	0.3317854	\\
40	0.3554958	\\
41	0.48654667	\\
42	0.49660314	\\
43	0.5230298	\\
44	0.75761317	\\
45	0.83516986	\\
46	0.95970209	\\
47	1.09713	\\
48	1.2602245	\\
49	1.2697649	\\
50	1.4709074	\\
51	1.4499263	\\
52	1.6115453	\\
53	1.7346132	\\
54	1.768639	\\
55	1.894712	\\
56	1.932978	\\
57	1.9691985	\\
58	2.1408199	\\
59	2.198014	\\
60	2.2522161	\\
61	2.329607	\\
62	2.4388831	\\
63	2.4980613	\\
64	2.5174607	\\
65	2.3850927	\\
66	2.403937	\\
67	2.3398796	\\
68	2.4810502	\\
69	2.4844838	\\
};
\addlegendentry{Data}

% Upper bound (standard error high)
\addplot[name path=upper, draw=none, forget plot] 
  table[row sep=crcr]{%,y expr=\thisrowno{1}*100]{%
30	0.31699318	\\
31	0.10011389	\\
32	0.35445438	\\
33	0.2506089	\\
34	0.29365724	\\
35	0.31921015	\\
36	0.36788288	\\
37	0.51851016	\\
38	0.51206232	\\
39	0.54050287	\\
40	0.55375253	\\
41	0.6729007	\\
42	0.67474134	\\
43	0.69266526	\\
44	0.9196122	\\
45	0.99146252	\\
46	1.1109516	\\
47	1.2443443	\\
48	1.4066153	\\
49	1.4124912	\\
50	1.6117372	\\
51	1.5897184	\\
52	1.74986	\\
53	1.8723206	\\
54	1.9050796	\\
55	2.0317431	\\
56	2.069126	\\
57	2.1103524	\\
58	2.2830432	\\
59	2.3452282	\\
60	2.3976214	\\
61	2.4770057	\\
62	2.5911442	\\
63	2.660554	\\
64	2.6842085	\\
65	2.563393	\\
66	2.5792043	\\
67	2.5160896	\\
68	2.66001	\\
69	2.6677378	\\
};

% Lower bound (standard error low)
\addplot[name path=lower, draw=none, forget plot] 
  table[row sep=crcr]{%,y expr=\thisrowno{1}*100]{%
30	-0.55490014	\\
31	-0.65290139	\\
32	-0.34352043	\\
33	-0.36422581	\\
34	-0.28304559	\\
35	-0.20369669	\\
36	-0.12917297	\\
37	0.04838584	\\
38	0.07000223	\\
39	0.12306792	\\
40	0.15724093	\\
41	0.30019264	\\
42	0.31846681	\\
43	0.35339433	\\
44	0.595616	\\
45	0.6788772	\\
46	0.80845444	\\
47	0.94991389	\\
48	1.1138318	\\
49	1.1270405	\\
50	1.3300777	\\
51	1.3101343	\\
52	1.4732306	\\
53	1.5969077	\\
54	1.6321967	\\
55	1.757679	\\
56	1.79683	\\
57	1.8280445	\\
58	1.9985966	\\
59	2.0507979	\\
60	2.1068127	\\
61	2.1822064	\\
62	2.2866238	\\
63	2.3355686	\\
64	2.3507148	\\
65	2.2067924	\\
66	2.2286696	\\
67	2.1636714	\\
68	2.3020905	\\
69	2.3012316	\\
};

% Fill between the paths
\addplot[blue!30] fill between[of=upper and lower];
\addlegendentry{95\% C.I.}

\addplot [color=red, line width=1.5pt, dashed]
  table[row sep=crcr]{%,y expr=\thisrowno{1}*100]{%
30	0	\\
31	0.16779459	\\
32	0.32122598	\\
33	0.4767707	\\
34	0.62655105	\\
35	0.76897051	\\
36	0.90586907	\\
37	1.0457174	\\
38	1.1805579	\\
39	1.3286312	\\
40	1.4692509	\\
41	1.6162862	\\
42	1.7495322	\\
43	1.8845477	\\
44	2.0202169	\\
45	2.155189	\\
46	2.2940789	\\
47	2.4300734	\\
48	2.5560681	\\
49	2.6679432	\\
50	2.777585	\\
51	2.8791545	\\
52	2.9829276	\\
53	3.0654487	\\
54	3.1348877	\\
55	3.1996158	\\
56	3.2695508	\\
57	3.3144833	\\
58	3.3506717	\\
59	3.3531731	\\
60	3.2686488	\\
61	3.2387078	\\
62	3.1352764	\\
63	3.1648646	\\
64	3.1213115	\\
65	3.0298607	\\
66	2.923937	\\
67	2.8263209	\\
68	2.7256512	\\
69	2.6072655	\\
  };
\addlegendentry{Model};

\end{axis}
\end{tikzpicture}
\end{tabular}}
\begin{minipage}{0.95\textwidth}\footnotesize\setstretch{1.0}
\textit{Notes:} Black: PSID profiles from \eqref{eq:profile}; shaded: 95 percent confidence bands; red dashed: model estimated with the utility function \eqref{eq:lnc}. Ages 30--69. (A) and (B) labor force participation in percent by health; (C) mean and (D) median nonpension assets in hundreds of thousands of 2016 dollars. ``Unhealthy'' pools poor health and disability. May 2025 model series. Source: PSID 1968--2015.
% src: F:h800_l0_lnc_edit.tex, h800_l1_lnc_edit.tex, h800_lnc_meana_paper.tex, h800_lnc_mediana_paper.tex (series styles, titles, axis labels; the _paper copies change the y-axis label only, figures/EDITS.md); 12:783-792 (May 2025 note)
\end{minipage}
\end{figure}
\FloatBarrier
